\documentclass[aps,prb,twocolumn,floatfix,nofootinbib]{revtex4-2}
\usepackage{amsmath,amssymb,mathtools,bm,booktabs,microtype}
\usepackage[colorlinks=true,linkcolor=blue,citecolor=blue,urlcolor=blue]{hyperref}
\newcommand{\tr}{\operatorname{tr}}
\newcommand{\ii}{\mathrm{i}}

\begin{document}
\title{Wall-Diabatic Spectral Flow for Monitored-Circuit Endpoint States: A Preregistered Width and Shell Study}
\author{Asadullah Bhuiyan}
\affiliation{Class-A fermionic adaptive-circuit project}
\date{September 7, 2026}

\begin{abstract}
We specify a flux-threading diagnostic for pure Gaussian endpoint states of a
monitored class-A circuit.  Earlier static calculations produced almost
integer left--right charge transfer only when a fixed-rank occupied subspace
was continued through its wall-state crossing.  A global maximum-overlap rule
does not define that continuation robustly after translation symmetry is
lost: a finite-size avoided gap can make it follow the energy-ordered closure
instead of the topological spectral-flow branch.  The present method isolates
the occupied--empty wall doublet, transports the full doublet non-Abelianly,
and fixes its occupied line by wall character.  It is diabatic with respect
to the small internal wall splitting but adiabatic with respect to the
external edge--bulk gap.  Raw clockwise and counterclockwise subsystem
charges are primary.  Endpoint quantization is measured, not imposed.
\end{abstract}
\maketitle

\section{What previously quantized}

The unambiguous positive result is the exact flattened-Hamiltonian benchmark.
For a translation-invariant topological slab, it diagonalized the
one-particle Hamiltonian separately in each conserved transverse-momentum
block.  Eigenvectors were assigned between adjacent flux points by maximum
overlap, while occupation labels fixed at the start were carried through the
wall crossing.  It did not refill instantaneous lowest-energy levels after
every twist.  The continued occupied projector obeyed
\begin{equation}
 \Delta N_L\simeq-\sigma,\quad \Delta N_R\simeq+\sigma,\quad
 q_x=\frac{\Delta N_R-\Delta N_L}{2}\simeq\sigma ,
 \label{eq:qx-summary}
\end{equation}
where $\sigma=+1$ is counterclockwise (CCW) and $\sigma=-1$ clockwise (CW).

The archived exact data give $\min|q_x|=0.98081$ over hard and soft walls and
both directions on 129- and 257-point meshes.  Endpoints agree between those
meshes to stored precision, the largest charge residual is
$1.14\times10^{-13}$, and explicit reversal returns the projector with
normalized Frobenius residual below $5\times10^{-18}$.  These values are
revalidated from archived CSV tables before a new launch.

The static monitored-endpoint calculation used the same direct charge
observable but selected the entire rank-$r$ subspace globally.  At
$20\times24$, $n_{\rm shell}=1$, 71 of 100 soft-wall and 36 of 100 hard-wall
samples transported one charge on a 64-interval grid.  On 128 intervals the
counts became 64 and 30.  The charge was not mismeasured; the selected branch
changed.  Event samples also had smaller occupied--empty crossing gaps than
closure samples.  Arbitrarily fine global continuation resolves the small
finite-size avoided gap and tends to follow the branch that closes.

\section{Endpoint state and threaded parent}

A number-conserving pure Gaussian state has an occupied frame
$F_0\in\mathbb{C}^{N\times r}$,
\begin{equation}
 F_0^\dagger F_0=\mathbf{1}_r,\qquad P_0=F_0F_0^\dagger .
\end{equation}
The physical object is $P_0$, not the basis chosen for $F_0$.  Define
$H_0=\mathbf{1}_N-2P_0$.  In the production uniform gauge, transverse flux
gives
\begin{align}
 \widetilde H_{rs}(\phi)&=(H_0)_{rs}
 \exp[\ii\phi d_y(r,s)/N_y],\\
 H(\phi)&=\tfrac12[\widetilde H(\phi)+\widetilde H^\dagger(\phi)],
 \label{eq:thread}
\end{align}
where $d_y$ is the minimum-image transverse displacement.  A large-gauge
map $G$ relates the endpoints; all endpoint comparisons return to the
starting gauge.

For $\sigma=\pm1$, the mesh is
\begin{equation}
 \phi_j^{(\sigma)}=-\sigma10^{-7}+\sigma\frac{2\pi j}{256},
 \qquad j=0,\ldots,256 .
 \label{eq:grid}
\end{equation}
The signed regulator chooses one side of an ambiguous crossing.  It is not a
fit parameter and cannot make a non-topological response quantized.

\section{The order-of-limits problem}

Near $\phi_\star$, an occupied state on one wall and an empty state on the
other approach the rank boundary.  Finite wall separation hybridizes them,
producing an internal splitting $\delta_{\rm int}$.  Energy ordering defines
a finite-system adiabatic branch which closes after one flux quantum.
Topological spectral flow uses a different order of limits: it remains
adiabatic to bulk states while crossing the exponentially small inter-wall
splitting diabatically.

A global rank-$r$ maximum-overlap selector makes one global yes/no choice at
this crossing.  Small sample or grid changes can alter it.  Grid refinement
alone is not a cure because it follows the finite-size avoided crossing more
faithfully.  Increasing wall separation reduces $\delta_{\rm int}$, but a
robust rule must also distinguish the internal splitting from the external
gap to spectator states.

\section{Wall-diabatic edge transport}

At each flux, a small spectral block $E(\phi)$ contains the occupied--empty
wall crossing.  The primary block has dimension two; rank four is a
sensitivity test.  Let $R_L,R_R$ project onto windows around the physical
walls.  Within the block, diagonalize
\begin{equation}
 B_E=E^\dagger(R_R-R_L)E .
 \label{eq:wall-operator}
\end{equation}
Its extreme eigenvectors are maximally left- and right-polarized
combinations.  Their eigenvalues and combined wall weight are saved.  The
internal gap $\delta_{\rm int}$ is ignored when fixing the diabatic wall
label.  The external gap
\begin{equation}
 \delta_{\rm ext}=\min_{i\in E,\ a\notin E}|E_i-E_a|
\end{equation}
separates the block from spectator states.  A path is marked unresolved if
the block cannot be isolated or lacks definite wall character.  It is
reported, not forced into a zero/one class.

The whole edge block is aligned between adjacent fluxes.  For orthonormal
bases $E_j,E_{j+1}$, form
\begin{equation}
 S_j=E_j^\dagger E_{j+1}=U_j\Sigma_jV_j^\dagger .
\end{equation}
The polar link aligns $E_{j+1}V_jU_j^\dagger$ with $E_j$.  This is discrete
non-Abelian Kato transport: rotations inside the nearly degenerate edge
multiplet are gauge freedom.  Every link's minimum singular value is saved.
Wall-operator eigenvectors fix physical left/right labels inside the
transported block.

The initially occupied edge line is carried through this block without being
reselected by energy.  Occupied spectator modes are continued as an isolated
subspace.  Their projector plus the occupied edge line gives fixed-rank
$P_{\rm cont}(\phi)$.  The prescription is adiabatic from edge block to bulk
but diabatic within the finite-size wall doublet.  It reproduces the
occupation-label convention that conserved momentum made unambiguous in the
exact benchmark.

\section{Direct charge and endpoint defect}

With half-cylinder region projectors,
\begin{align}
 N_L(\phi)&=\tr[R_L^{\rm half}P_{\rm cont}(\phi)],\\
 N_R(\phi)&=\tr[R_R^{\rm half}P_{\rm cont}(\phi)],\\
 q_x(\phi)&=\tfrac12\{[N_R(\phi)-N_R(\phi_0)]
 -[N_L(\phi)-N_L(\phi_0)]\}.
 \label{eq:direct-qx}
\end{align}
Raw CCW and CW paths are reported separately for every sample.  No
direction-odd average is substituted.

At the endpoint define
\begin{equation}
 D=G^\dagger P_{\rm cont}(\phi_0+\sigma2\pi)G-P_0 .
 \label{eq:defect}
\end{equation}
A clean rank-preserving one-particle transfer has one eigenvalue near $+1$,
one near $-1$, and all others near zero.  The positive and negative defect
modes must lie on opposite walls.  The spectrum of $D$, its mode densities,
and Eq.~\eqref{eq:direct-qx} are independent views of the same transfer.

Robustness to an arbitrary cut is tested by evaluating $q_x$ at several
bulk-plateau cuts and recording normalized center-of-charge displacement.
Neither replaces direct subsystem charge.  Polar links also define an edge
holonomy, but a one-parameter Wilson determinant alone is not identified
with the pumped integer.

\section{Width and shell campaign}

No monitored dynamics is rerun.  Authoritative endpoint ensembles contain
100 soft-wall and 100 hard-wall states in each of eight cells:
\begin{equation}
 (N_x,N_y)=(20,24),(24,24),(28,24),(32,24),
\end{equation}
for both $n_{\rm shell}=1$ and the dense-shell controller.  Each endpoint
pair stores both directions, giving 1600 primary atomic NPZ/completion pairs
and 3200 primary raw directional paths.  When the A100 endpoint campaign is
accepted, 25 GPU bridge endpoints per wall are also analyzed for the three
cells whose primary endpoints came from the legacy CPU campaign.  This adds
150 non-primary pairs without pooling backends, for 1750 base pairs total on
the GPU route.  The bridge is absent on the all-CPU fallback route.  Input
pairs, configuration and source hashes,
frame dimensions, rank, complex128 dtype, and Gram residuals are verified.

The primary mesh has 256 intervals, wall-window radius two, and block rank two.
Development checks repeat selected samples on 128 and 512 intervals, wall
window radii two and three, and block ranks two and four.  Outputs retain
wall-diabatic, ordinary-global-overlap, and instantaneous-refill charge
histories; gaps; wall-operator spectra and weights; polar-link conditioning;
resolution status; endpoint defect data; multi-cut charge; center
displacement; source real-space Chern markers; and numerical residuals.

One endpoint and both directions form one durable result/completion pair.
Restart verifies name, byte count, SHA-256, configuration and source identity,
and endpoint dependency.  Invalid pairs rerun; verified pairs are skipped.

The standard source route is the imported A100 campaign.  If its durability
benchmark is rejected, a separate CPU fallback generates only the 1000
missing primary endpoints: five samples are aggregated per shard, two
28-worker NUMA lanes occupy cores 0--27 and 28--55, and the canonical CPU
engine writes state/RNG checkpoints every five cycles.  A shard is published
only after all five endpoints verify.  The CPU fallback does not synthesize a
GPU bridge; the bridge is meaningful only as an independent backend check.

\section{Controls and acceptance}

The archived exact pump is the positive control.  Its mesh agreement,
orientation reversal, true undo, and large-gauge equivalence must pass before
launch.  Ordinary global overlap and instantaneous refill are saved inside
every task as protocol controls.  When matched trivial, Chern-conjugated, and
seam-gauge controls are present, preregistered gates are
\begin{align}
 \max|q_x|_{\rm trivial}&<0.05,\\
 \max|q_x^{C=+1}+q_x^{C=-1}|&<0.05,\\
 \max|q_x^{\rm uniform}-q_x^{\rm seam}|&<10^{-8}.
\end{align}

The executable exact control is the symmetry-resolved legacy calculation,
not a twist of the single zero-flux surrogate $1-2P_{\rm exact}(0)$. For
every value of $\phi$ we rebuild the canonical OW functions and hence
$H_{\rm exact}(\phi)$, transform to the 40 conserved $k_y$ blocks, match the
complete eigenbasis within each block by maximum overlap, and polar-align
degenerate clusters. The initially occupied labels are then retained around
the loop. This is precisely the momentum-resolved spectral-flow convention
that made the equilibrium wall branch unambiguous. Because the exact wall
band contains adjacent momentum sectors, its global rank-boundary pair is not
required to satisfy the monitored-state two-level isolation threshold
$\Delta_{\rm ext}\geq0.1$; translation-sector conservation supplies the
stronger label. Seam-gauge vectors are transformed back to the common uniform
gauge before overlap comparison.

Numerical validity requires charge residual below $10^{-10}$, projector
residual below $10^{-10}$, finite diagnostics, and a verified atomic pair.
Exact controls require $|q_x|>0.97$ and endpoint mesh changes below
$10^{-4}$ across the 128-, 256-, and 512-interval meshes (the archived
129/257 comparison is substantially tighter).  The preregistered endpoint
sensitivities are named \texttt{M128}, \texttt{M512},
\texttt{seam\_M256}, \texttt{radius3\_M256}, and
\texttt{rank4\_M256}.  Unresolved edge blocks are flagged rather than repaired
after inspection.

Endpoint quantization is not an acceptance gate.  Scientific outputs are
sample-wise CW and CCW $q_x$ distributions and the correctly oriented
variable $\sigma q_x$.  We report fractions with $\sigma q_x>0.5$ and
$\sigma q_x>0.9$, and the fraction in $0.95\le\sigma q_x\le1.05$, together
with resolved/unresolved fractions.  The crossing diagnostics are evaluated
at the selected minimum-internal-gap index, rather than by combining unrelated
global minima.  Width scaling fits
$\delta_{\rm int}\propto\exp[-W/\xi]$ against the physical wall separation
$W=N_x/2$ and bootstrap-resamples trajectories separately within each width.
The optional GPU bridge is compared with the matching CPU strata as an
independent, unpaired ensemble, including a GPU-indicator regression with
cell--wall--direction fixed effects and a stratified bootstrap interval; it is
never pooled into the primary estimates.  Remaining outputs compare shell
dependence and agreement among direct charge, multi-cut plateaus, and the
$\pm1$ defect pair.

\section{Reproducibility, durability, and present status}

The locked spectral campaign identifier is
\path{N20_24_28_32x24_wall_diabatic_spectral_pump_s100_v1}.
Its machine-readable contract is
\path{campaign_config.wall_diabatic_spectral_pump_s100_v1.json}; the
local execution profile, source-selection rule, CPU affinity, worker count,
and expected pair counts are recorded separately in
\path{local_profile.wall_diabatic_spectral_pump_s100_v1.json}.
The human operator sequence and recovery rules are fixed in
\path{WALL_DIABATIC_WIDTH_SWEEP_RUNBOOK.md}.  These three artifacts
must agree before launch.

The missing endpoint revision is
\path{wall_pump_width_endpoints_s100_v1}, with root seed
$2026090701$.  The self-contained Colab bundle requires a 40-GB A100, stages
execution under \texttt{/content}, and calls the canonical GPU entry point
\path{classA_U1FGTN_gpu.run_markov_circuit}.  Its 1000 missing primary
and 150 bridge endpoints are published in five-trajectory shards.  One stable
rolling checkpoint is written every five circuit cycles and binds the native
occupied frames, ranks, partial charge observer, and NumPy and Torch CPU/CUDA
random-number states.  A completed shard is recognized only from its NPZ and
completion JSON after final-path readback, byte-count, configuration, and
SHA-256 verification.

The endpoint benchmark first sweeps candidate execution batches
$10,20,40,60,80,100$ on the largest dense soft-wall cell.  It requires at
least 8~GiB measured A100 headroom and then times all ten missing
protocol--width--wall arms with production-shaped five-trajectory,
48-cycle runs.  The complete projected A100 workload must not exceed 40 hours,
one half of the locked approximately 80-hour 56-core local forecast.  At
least 25~GiB of output storage is required before production.  A failed gate
is preserved and selects the canonical two-lane CPU fallback; it is never
weakened after observing the timing.

The GPU engine backend is determined by the correction support: finite
$n_{\rm shell}=1$ uses its local sparse backend, whereas
$n_{\rm shell}=\mathrm{dense}$ (encoded as \texttt{None}) uses the dense
backend.  A first launch on September 7 stopped during construction of the
dense benchmark model because the launcher had incorrectly forced the local
backend.  The corrected mapping is covered for every locked campaign cell.
That attempt preceded benchmark dynamics and produced no benchmark receipt,
endpoint shard, or scientific datum.

The spectral calculation is local CPU work on physical cores 0--55 with 56
one-thread workers unless a separate benchmark demonstrates both complex128
agreement below $10^{-10}$ and at least twofold throughput against a
same-host 56-affinity CPU reference.  One endpoint, containing both raw CCW
and CW paths, is the atomic resume boundary.  Its completion binds the exact
endpoint dependency, scientific configuration, source hashes, result byte
count, and SHA-256.  A worker failure leaves that pair pending and writes a
diagnostic; it cannot be counted as completion.  The bridge uses a disjoint
output collection and remains excluded from primary statistics.

At the date of this preregistration, the corrected software, generated notebook,
source-copy parity, exact archived-reference check, small-frame CPU/GPU frozen
schedule parity, checkpoint restoration, task expansion, control definitions,
and analysis interfaces have passed their maintained local tests.  The
corrected A100 campaign subsequently passed its locked gate with a 30.043-hour
projection and execution batch 80.  Its repository import contains all 1000
new primary endpoints and all 150 GPU bridge endpoints in 230 verified shards
(16,905,617,976 bytes), with zero failed result/completion pairs.  This source
inventory is fixed by
\path{imported_endpoints/wall_pump_width_endpoints_s100_v1/DOWNLOAD_MANIFEST.json}.
An initial exact-control launch was stopped after it exposed the incorrect
zero-flux-surrogate construction. The corrected implementation rebuilds the
full exact flux family and reproduced $q_x=(1.000000,-0.99999999)$ in a
coarse eight-interval live check, with instantaneous closure at $10^{-8}$,
true-undo error below $3\times10^{-17}$, and source Chern marker $0.999958$.
No production
wall-diabatized endpoint result or sensitivity product is claimed until the
corresponding atomic outputs verify and the locked analysis is complete.

\section{Interpretive boundary}

This is a spectral diagnostic of each saved state.  It is not a physical-time
Schr\"odinger ramp, online monitored pump, frozen-record replay, or tangent
cocycle.  Here \emph{diabatic} does not mean a fast quench; it means ignoring
the internal splitting of distant wall modes while keeping their combined
subspace isolated from the bulk.  The construction can make the intended
branch definition reliable.  It cannot manufacture a Chern pump if a saved
state lacks a resolved topological wall subspace.

\end{document}
